\chapter{Simulation Output Metrics Reference}
\label{appendix:metrics}

This appendix documents every metric produced by the co-simulation framework
and displayed in the Streamlit dashboard (\texttt{app.py}).  For each scenario
the metrics are grouped by \emph{domain concern} and described in terms of
what they measure, how they are computed in the simulation code, and what
simulation-level event or physical process they reflect.  Equation and section
cross-references point back to the corresponding models in
Chapter~\ref{chapter:methodology}.

Each entry follows the format:
\begin{quote}
\texttt{metric\_key} \quad [\textit{unit}] --- description.
\end{quote}

% ------------------------------------------------------------------
\section{Scenario E1 --- Smart Grid with Renewable Integration}
\label{sec:metrics-e1}

\subsection*{Energy Generation}

\begin{description}

\item[\texttt{total\_renewable\_generation\_kwh}] [kWh] ---
  Total electrical energy produced by all solar PV and wind turbines over the
  24-hour window.  Computed by summing instantaneous combined output at every
  15-minute step and multiplying by $\Delta t = 0.25\,\text{h}$.  Compare
  with \texttt{total\_load\_kwh} to assess net self-sufficiency.

\item[\texttt{total\_solar\_generation\_kwh}] [kWh] ---
  Energy contribution from solar PV units alone.  Each unit follows the
  sine-based irradiance model (Equation~\ref{eq:solar_power}), producing
  output only during daylight hours (06:00--18:00).  Stochastic cloud factors
  introduce step-to-step variability.

\item[\texttt{total\_wind\_generation\_kwh}] [kWh] ---
  Energy contribution from wind turbines alone.  Each turbine follows the
  three-region power curve (Equation~\ref{eq:wind_power}), with zero output
  below cut-in ($3\,\text{m/s}$) and above cut-out ($25\,\text{m/s}$).

\item[\texttt{renewable\_penetration\_pct}] [\%] ---
  Fraction of total demand served by renewables:
  $\text{penetration} = P_{renewable} / P_{load} \times 100$.  Values above
  100\% indicate periods of surplus generation requiring battery storage or
  curtailment.

\end{description}

\subsection*{Load and Demand Response}

\begin{description}

\item[\texttt{total\_load\_kwh}] [kWh] ---
  Aggregate energy consumed by all residential loads over 24 hours.  Each of
  the 800 \texttt{ResidentialLoad} agents uses a time-varying load factor
  (0.5 off-peak to 1.0 during evening peak 17:00--22:00) with $\pm10\%$
  stochastic variation.

\item[\texttt{peak\_load\_kw}] [kW] ---
  Maximum instantaneous aggregate demand across all 96 time steps.  Determines
  worst-case feeder loading and indicates whether grid reinforcement would be
  required in a real deployment.

\item[\texttt{avg\_peak\_hour\_load\_kw}] [kW] ---
  Mean load during evening peak hours (17:00--22:00) after demand response is
  applied.  Useful for evaluating how effectively DR flattens the evening peak.

\item[\texttt{avg\_off\_peak\_load\_kw}] [kW] ---
  Mean load outside peak hours; the baseline against which peak shaving is
  measured.

\item[\texttt{estimated\_dr\_load\_reduction\_kw}] [kW] ---
  Estimated load reduction attributable to demand response: difference between
  the expected peak-hour load without DR and the actual load with DR applied.
  DR activates for the 70\% of loads that opt in
  (\texttt{dr\_participation = True}), cutting consumption by 15\% during
  on-peak hours (Section~\ref{sec:scenario-e1}).

\item[\texttt{estimated\_dr\_load\_reduction\_pct}] [\%] ---
  Same reduction expressed as a percentage of the pre-DR peak load.

\end{description}

\subsection*{Battery Storage}

\begin{description}

\item[\texttt{avg\_battery\_soc}] [---] ---
  Time-averaged state of charge across all BESS units, normalised to $[0,1]$.
  Updated each step via Equation~\ref{eq:soc_charge} with
  $\eta_{charge} = 0.92$.  Tracks how actively batteries buffer renewable
  surplus or cover deficits across the day.

\item[\texttt{final\_battery\_soc}] [---] ---
  List of per-unit SOC values at the end of the 24-hour simulation.  Displayed
  as a bar chart in the dashboard.  Units close to $SOC_{max} = 0.95$ indicate
  incomplete discharge; units near $SOC_{min} = 0.20$ indicate heavy use.

\end{description}

\subsection*{Curtailment and Voltage}

\begin{description}

\item[\texttt{total\_curtailment\_kwh}] [kWh] ---
  Renewable energy discarded because batteries were fully charged and no
  additional load was available.  Curtailment occurs when
  $P_{net}^{final}(t) < 0$ even after maximum battery charging
  (Section~\ref{sec:scenario-e1}).

\item[\texttt{curtailment\_percentage}] [\%] ---
  Curtailment as a fraction of total renewable generation.  High values
  suggest insufficient storage or excess installed capacity relative to demand.

\item[\texttt{curtailment\_steps}] [steps] ---
  Number of 15-minute intervals in which any curtailment occurred.

\item[\texttt{curtailment\_with\_saturated\_battery}] [steps] ---
  Subset of curtailment steps where at least one battery reached $SOC_{max}$.
  Used as the denominator for the validation check below.

\item[\texttt{curtailment\_battery\_validation}] [PASS/FAIL] ---
  Internal consistency check: curtailment may only occur when batteries are
  saturated.  A \texttt{FAIL} indicates a logic error in
  \texttt{calculate\_storage\_dispatch()}.

\item[\texttt{avg\_voltage\_pu}] [p.u.] ---
  Mean bus voltage across all 33 IEEE buses and all time steps.  Computed via
  the simplified linear drop model
  $V_{bus}(i) = V_{base} - k_{drop} \cdot P_{bus}(i) \cdot i$.
  ANSI~C84.1 requires $V \in [0.95,\,1.05]$\,p.u.

\item[\texttt{min\_voltage\_pu}] [p.u.] ---
  Minimum bus voltage recorded.  Values below 0.95 signal under-voltage events
  on the most heavily loaded buses.

\item[\texttt{max\_voltage\_pu}] [p.u.] ---
  Maximum bus voltage recorded.  Values above 1.05 indicate over-voltage from
  excessive reverse power flow during surplus generation.

\item[\texttt{voltage\_violations}] [count] ---
  Number of time steps in which any bus voltage fell outside the ANSI limits.
  Non-zero values would require voltage regulation in a real grid.

\end{description}

% ------------------------------------------------------------------
\section{Scenario E2 --- Electric Vehicle Charging Infrastructure}
\label{sec:metrics-e2}

\subsection*{Grid Impact}

\begin{description}

\item[\texttt{peak\_load\_kw}] [kW] ---
  Highest aggregate load (base load plus EV charging) recorded during the
  36-hour window.  The uncoordinated strategy typically peaks during the
  residential evening (17:00--22:00) when EV arrivals and household demand
  coincide.  Smart and V2G strategies spread and suppress this peak.

\item[\texttt{max\_transformer\_loading\_pct}] [\%] ---
  Highest per-transformer loading fraction relative to rated capacity.  Values
  above 100\% indicate the transformer was overloaded and would trip protection
  equipment in a real network.  Overload events are counted separately.

\item[\texttt{avg\_transformer\_loading\_pct}] [\%] ---
  Mean transformer utilisation over time; a proxy for infrastructure wear.

\item[\texttt{transformer\_overload\_count}] [count] ---
  Number of time steps in which any transformer exceeded rated capacity.
  Directly linked to \texttt{peak\_load\_kw}: peaks above feeder capacity
  generate overloads.

\item[\texttt{max\_voltage\_deviation\_pu}] [p.u.] ---
  Maximum deviation of any node voltage from the nominal 1.0\,p.u., used as a
  proxy for feeder voltage quality.  V2G power injections during on-peak hours
  can cause positive deviations (over-voltage).

\item[\texttt{load\_factor}] [---] ---
  Ratio of average to peak load over the simulation horizon.  Values close to
  1.0 indicate a flat, well-managed profile; values near 0 indicate sharp
  peaks.  Smart and V2G charging raise the load factor by shifting demand to
  off-peak hours.

\end{description}

\subsection*{Charging Economics and EV Outcomes}

\begin{description}

\item[\texttt{total\_energy\_charged\_kwh}] [kWh] ---
  Total electrical energy delivered to EV batteries.  Each charge step uses the
  urgency-based allocation (Equation~\ref{eq:urgency}) for smart/V2G modes, or
  immediate full-rate delivery for uncoordinated.

\item[\texttt{total\_charging\_cost\_usd}] [USD] ---
  Aggregate cost paid by all EV owners under the TOU tariff
  (Equation~\ref{eq:tou}).  Smart charging preferentially uses off-peak
  periods (0.08\,USD/kWh) to minimise cost.

\item[\texttt{avg\_cost\_per\_vehicle\_usd}] [USD] ---
  \texttt{total\_charging\_cost\_usd} divided by the number of vehicles.
  Normalises fleet-size effects when comparing strategies.

\item[\texttt{avg\_final\_soc}] [---] ---
  Fleet-average state of charge at simulation end.  A value below
  $SOC_{target} = 0.80$ indicates that some vehicles could not fully charge
  due to grid or time constraints.

\item[\texttt{vehicles\_connected}] [count] ---
  Number of vehicles that successfully connected to a charging station during
  the simulation.

\item[\texttt{vehicles\_meeting\_soc\_target}] [count] ---
  Vehicles that reached $SOC_{target} = 0.80$ before departure.  The primary
  user-satisfaction metric for the charging infrastructure.

\item[\texttt{v2g\_energy\_provided\_kwh}] [kWh] ---
  Energy exported from EV batteries back to the grid during on-peak hours
  (17:00--21:00) by V2G-capable vehicles
  (50\% of fleet, $P_{discharge}^{max} = 10\,\text{kW}$).

\item[\texttt{v2g\_revenue\_usd}] [USD] ---
  Revenue earned from grid export, compensated at 20\% above the current TOU
  rate.  Non-zero only under the V2G strategy.

\end{description}

% ------------------------------------------------------------------
\section{Scenario M1 --- Urban Traffic Congestion Management}
\label{sec:metrics-m1}

\subsection*{Vehicle Progress}

\begin{description}

\item[\texttt{completed\_vehicles}] [count] ---
  Vehicles that reached their destination within the 3-hour window.  The 2,500
  agents depart stochastically during the first hour; completion depends on A*
  route length and signal-induced delays.

\item[\texttt{completion\_rate\_pct}] [\%] ---
  Percentage of the fleet that completed their trip.  A lower rate under
  fixed-time signalling compared to adaptive signals quantifies the throughput
  benefit of adaptive control.

\item[\texttt{avg\_travel\_time\_s}] [s] ---
  Mean door-to-door travel time for completed vehicles.  Free-flow travel time
  is $\approx 90\,\text{s/link} \times$ path length; excess time reflects
  signal waiting and queue delays.

\item[\texttt{avg\_travel\_time\_min}] [min] ---
  Same metric in minutes for readability in the dashboard.

\item[\texttt{avg\_delay\_s}] [s] ---
  Difference between actual and free-flow travel time, averaged over completed
  vehicles.  The primary efficiency metric contrasting adaptive and fixed
  signal strategies.

\item[\texttt{avg\_delay\_min}] [min] ---
  Same metric in minutes.

\item[\texttt{total\_delay\_s}] [s] ---
  Sum of individual delays across all vehicles.  Scales with fleet size;
  normalise by \texttt{completed\_vehicles} when comparing runs with different
  fleet sizes.

\item[\texttt{throughput\_veh\_per\_hour}] [veh/h] ---
  Rate of trip completions: \texttt{completed\_vehicles} divided by simulation
  duration in hours.  Higher values indicate better network efficiency.

\end{description}

\subsection*{Queues and Emissions}

\begin{description}

\item[\texttt{max\_queue\_length}] [vehicles] ---
  Longest queue observed at any intersection approach during the simulation.
  Reflects worst-case congestion; under fixed signals it is typically higher
  because green time is not directed toward the most loaded approach.

\item[\texttt{avg\_queue\_length}] [vehicles] ---
  Mean queue length averaged across all intersections and all time steps.  A
  proxy for overall network saturation.

\item[\texttt{total\_emissions\_kg\_co2}] [kg\,CO$_2$] ---
  Total CO$_2$ produced by the entire fleet, computed using the two-state model
  (Equation~\ref{eq:emissions}): idle vehicles emit
  $\alpha_{idle} = 2.31\,\text{g/s}$; moving vehicles emit
  $\alpha_{move} = 0.15\,\text{g/m}$.  Longer red-light waits directly raise
  idle emissions.

\item[\texttt{emissions\_per\_vehicle\_g}] [g\,CO$_2$] ---
  Mean emissions per vehicle in grams.  Normalises total emissions by fleet
  size to allow fair comparison between differently sized runs.

\end{description}

% ------------------------------------------------------------------
\section{Scenario T1 --- 5G Slice Resource Allocation}
\label{sec:metrics-t1}

\subsection*{QoS and Slice Utilisation}

\begin{description}

\item[\texttt{qos\_satisfaction\_embb\_pct}] [\%] ---
  Fraction of eMBB user-steps in which the user received at least
  $RB_{req} = 2$ resource blocks (Equation~\ref{eq:qos_t1}).  eMBB users
  require high throughput for broadband services; satisfaction drops when their
  slice allocation is reduced under heavy demand from other slices.

\item[\texttt{qos\_satisfaction\_urllc\_pct}] [\%] ---
  Fraction of URLLC user-steps meeting the $RB_{req} = 2$ threshold.  URLLC
  users require ultra-reliable, low-latency connectivity; failed handovers
  automatically count as QoS failures.  This metric drives the coupling
  feedback path in Scenario~M1+T1 (Section~\ref{sec:scenario-m1t1}).

\item[\texttt{qos\_satisfaction\_mmtc\_pct}] [\%] ---
  Fraction of mMTC device-steps meeting $RB_{req} = 1$.  mMTC devices follow a
  bursty activity model with base probability 30\% and 50-second burst windows
  (Section~\ref{sec:scenario-t1}).

\item[\texttt{slice\_utilization\_embb\_pct}] [\%] ---
  Fraction of RBs allocated to the eMBB slice actually consumed by active
  users.  Static slicing often shows low utilisation for one slice while
  another is starved.

\item[\texttt{slice\_utilization\_urllc\_pct}] [\%] ---
  Same as above for the URLLC slice.

\item[\texttt{slice\_utilization\_mmtc\_pct}] [\%] ---
  Same as above for the mMTC slice.

\item[\texttt{overall\_utilization\_pct}] [\%] ---
  Fraction of total available RBs (across all gNBs and slices) that are
  consumed.  High values indicate a well-loaded network; low values suggest
  over-provisioning or poor load balancing.

\item[\texttt{resource\_waste\_pct}] [\%] ---
  RBs allocated to a slice but not used by any active user in a given step,
  as a fraction of total available RBs.  Static slicing typically produces
  higher waste because fixed quotas do not track instantaneous demand.

\end{description}

\subsection*{Handovers}

\begin{description}

\item[\texttt{handover\_attempts}] [count] ---
  Total handover events triggered during the simulation.  A handover occurs
  when a neighbouring gNB's received signal exceeds the serving gNB's by the
  3\,dB hysteresis margin, subject to a 10-step cooldown to suppress
  ping-pong effects (Section~\ref{sec:scenario-t1}).

\item[\texttt{handover\_success\_rate\_pct}] [\%] ---
  Fraction of attempted handovers that succeeded.  Each attempt succeeds with
  probability $P_{success} = 0.95$; a failure causes one step of QoS outage
  for that user.

\item[\texttt{handover\_failures}] [count] ---
  Number of unsuccessful handovers.  Each failure is also counted as one QoS
  failure in the affected user's slice, directly reducing the corresponding
  \texttt{qos\_satisfaction\_*\_pct} value.

\end{description}

% ------------------------------------------------------------------
\section{Scenario E2+M1 --- Cross-Domain Energy--Mobility Integration}
\label{sec:metrics-e2m1}

The E2+M1 scenario produces metrics from both domain modules plus
coupling-specific metrics that are only meaningful when the two domains
interact.

\subsection*{Background Traffic}

These metrics characterise the 2,000 background vehicles that share the road
grid with the EV fleet and are affected by the congestion EVs create.

\begin{description}

\item[\texttt{bg\_completion\_rate\_pct}] [\%] ---
  Fraction of background vehicles completing their trip.  In coupled mode, EVs
  consume road capacity and cause signal queues, reducing throughput for
  background traffic compared to uncoupled mode.

\item[\texttt{bg\_avg\_travel\_time\_s}] [s] ---
  Mean travel time for background vehicles.  The difference between coupled and
  uncoupled values isolates the road-capacity impact of EV navigation.

\item[\texttt{bg\_avg\_delay\_s}] [s] ---
  Mean additional delay beyond free-flow travel time for background vehicles.
  Increases in coupled mode as EVs add to intersection queues.

\item[\texttt{max\_queue\_length}] [vehicles] ---
  Worst-case queue at any intersection approach, including both EV and
  background vehicles.

\item[\texttt{avg\_queue\_length}] [vehicles] ---
  Network-average queue length including all vehicle types.

\item[\texttt{total\_emissions\_kg\_co2}] [kg\,CO$_2$] ---
  Total CO$_2$ from all vehicles (EVs and background) via the two-state model
  (Equation~\ref{eq:emissions}).  EVs produce emissions while driving to
  stations; their idle time at red lights constitutes the \emph{emission
  uplift} coupling metric (Section~\ref{sec:scenario-e2m1}).

\end{description}

\subsection*{EV Travel and Coupling}

\begin{description}

\item[\texttt{avg\_ev\_drive\_time\_s}] [s] ---
  Mean time an EV spends actively moving toward its assigned charging station
  (phase \texttt{DRIVING\_TO\_STATION}, Equation~\ref{eq:ev_phases}).
  Approximately 0\,s in uncoupled mode (teleportation); non-zero in coupled
  mode due to traffic delays.

\item[\texttt{avg\_ev\_queue\_time\_s}] [s] ---
  Mean time an EV spends waiting at red signals during its station-bound trip.
  The primary coupling signal: queuing time is exactly zero in uncoupled mode.

\item[\texttt{avg\_ev\_travel\_to\_station\_s}] [s] ---
  End-to-end elapsed time from EV departure to first arrival at the charging
  station intersection node.

\item[\texttt{evs\_at\_station\_node}] [count] ---
  Number of EVs that reached the intersection containing their assigned station
  (but may not yet have connected to a port).

\item[\texttt{evs\_reached\_station}] [count] ---
  Number of EVs that physically connected to a charging port.

\item[\texttt{evs\_meeting\_soc\_target}] [count] ---
  Number of EVs that reached $SOC_{target}$ before simulation end.
  Travel-delayed arrivals consume available charging time, so coupled mode
  typically records fewer vehicles meeting the target; this gap is the
  \emph{SOC target deficit} coupling metric
  (Section~\ref{sec:scenario-e2m1}).

\item[\texttt{avg\_final\_soc}] [---] ---
  Fleet-average EV state of charge at simulation end.

\end{description}

\subsection*{Charging and Grid}

\begin{description}

\item[\texttt{total\_energy\_charged\_kwh}] [kWh] ---
  Total energy delivered to all EVs that reached a station.  Reduced in
  coupled mode relative to uncoupled because travel delays shrink the
  available charging window.

\item[\texttt{total\_charging\_cost\_usd}] [USD] ---
  Aggregate cost under the TOU tariff.  Coupled arrivals at later hours may
  fall into on-peak windows, increasing total cost.

\item[\texttt{peak\_grid\_load\_kw}] [kW] ---
  Maximum concurrent power draw from the charging infrastructure.  DC fast
  stations (50\,kW each) at inner grid intersections dominate the peak.

\item[\texttt{max\_transformer\_loading\_pct}] [\%] ---
  Highest observed transformer utilisation.  Same interpretation as in
  Scenario~E2 (Section~\ref{sec:metrics-e2}).

\item[\texttt{transformer\_overloads}] [count] ---
  Number of 5-minute charging steps in which any transformer exceeded rated
  capacity.

\item[\texttt{v2g\_energy\_kwh}] [kWh] ---
  Energy exported to the grid by V2G-capable EVs.  Non-zero only when the V2G
  charging strategy is selected.

\item[\texttt{v2g\_revenue\_usd}] [USD] ---
  Revenue from V2G exports at the 20\% on-peak premium rate.

\end{description}

% ------------------------------------------------------------------
\section{Scenario M1+T1 --- Cross-Domain Mobility--Telecommunications Integration}
\label{sec:metrics-m1t1}

The M1+T1 scenario produces metrics from both domain modules plus
coupling-specific metrics that expose the causal chain from vehicular
clustering through URLLC QoS degradation to traffic signal impairment.

\subsection*{Telecom QoS and Resource Allocation}

\begin{description}

\item[\texttt{qos\_satisfaction\_embb\_pct}] [\%] ---
  Fraction of eMBB user-steps meeting $RB_{req} = 2$
  (Equation~\ref{eq:qos_t1}).  In coupled mode, vehicular clustering
  concentrates users near cell boundaries, reducing per-user RB allocation
  compared to the spatially uniform RWP baseline.

\item[\texttt{qos\_satisfaction\_urllc\_pct}] [\%] ---
  Critical metric driving the QoS feedback path.  When this drops below
  $\theta_{URLLC} = 0.80$ at a gNB, all intersections within that gNB's
  coverage radius revert to fixed-time signal control
  (Section~\ref{sec:scenario-m1t1}).

\item[\texttt{qos\_satisfaction\_mmtc\_pct}] [\%] ---
  Fraction of mMTC device-steps meeting $RB_{req} = 1$.  mMTC devices are
  intermittent (30\% base activity), so this metric is less sensitive to
  congestion than eMBB or URLLC.

\item[\texttt{overall\_utilization\_pct}] [\%] ---
  Fraction of all available RBs consumed.  Vehicular clustering can raise
  utilisation at affected gNBs while leaving others under-utilised, explaining
  the elevated per-gNB variance in coupled mode.

\item[\texttt{resource\_waste\_pct}] [\%] ---
  Unused allocated RBs as a fraction of total.  Static slicing produces more
  waste than dynamic in both coupled and uncoupled modes.

\end{description}

\subsection*{Handovers and Load Imbalance}

\begin{description}

\item[\texttt{handover\_attempts}] [count] ---
  Total handover triggers across all connected vehicles.  Vehicles travelling
  at 50\,km/h produce rapid cell transitions, whereas RWP pedestrians
  (0.5--2.0\,m/s) cross boundaries slowly.  Coupled mode is expected to show
  substantially more attempts.

\item[\texttt{handover\_successes}] [count] ---
  Successfully completed handovers (success probability 0.95).

\item[\texttt{handover\_failures}] [count] ---
  Failed handovers, each causing one step of QoS outage for the affected user
  and one failure counted in the corresponding QoS satisfaction metric.

\item[\texttt{handover\_success\_rate\_pct}] [\%] ---
  \texttt{handover\_successes} / \texttt{handover\_attempts} $\times 100$.
  A lower value in coupled mode reflects the higher volume of bursty handovers
  from stop-and-go vehicle motion.

\item[\texttt{handover\_rate\_per\_min}] [events/min] ---
  Handover attempts normalised by simulation duration.  The difference between
  coupled and uncoupled values is the \emph{handover rate uplift} coupling
  metric (Section~\ref{sec:scenario-m1t1}).

\item[\texttt{load\_imbalance\_std\_users}] [users] ---
  Standard deviation of active users per gNB, averaged over all simulation
  steps.  Vehicular clustering at intersections near gNB coverage boundaries
  pushes users unevenly across cells, raising this value above the near-zero
  RWP baseline.

\end{description}

\subsection*{Signal Degradation and Traffic}

\begin{description}

\item[\texttt{degradation\_events}] [count] ---
  Number of distinct gNB degradation activations: each time a gNB's URLLC
  satisfaction rate falls below 0.80 and one or more of its covered
  intersections receive the \texttt{signal\_degraded} flag.

\item[\texttt{degradation\_duration\_s}] [s] ---
  Cumulative seconds during which adaptive signal control was suppressed at one
  or more intersections.  This is the \emph{signal degradation duration}
  coupling metric (Section~\ref{sec:scenario-m1t1}).

\item[\texttt{affected\_intersections}] [count] ---
  Total number of distinct intersection nodes that received the
  \texttt{signal\_degraded} flag at least once.  Always zero in uncoupled mode.

\item[\texttt{bg\_completion\_rate\_pct}] [\%] ---
  Fraction of background vehicles completing their trip.  Degrades when
  adaptive signals are suppressed, reducing intersection throughput.

\item[\texttt{bg\_avg\_travel\_time\_s}] [s] ---
  Mean travel time of background (non-connected) vehicles.  Increases during
  degradation periods because affected intersections revert to fixed
  equal-phase timing ($Q_{NS} = Q_{EW} = 1$,
  Equation~\ref{eq:adaptive_signal}).

\item[\texttt{bg\_avg\_delay\_s}] [s] ---
  Mean excess delay beyond free-flow for background vehicles.  The difference
  between coupled and uncoupled values quantifies the traffic cost of the QoS
  feedback path.

\item[\texttt{cv\_avg\_travel\_time\_s}] [s] ---
  Mean travel time of the 200 connected vehicles.  These vehicles are doubly
  affected: they contribute to the radio congestion that triggers signal
  degradation, and they also experience the resulting traffic delays.

\item[\texttt{total\_emissions\_kg\_co2}] [kg\,CO$_2$] ---
  Total CO$_2$ from all vehicles.  Signal degradation increases idle time at
  affected intersections, raising per-vehicle emissions via
  $\alpha_{idle} = 2.31\,\text{g/s}$ (Equation~\ref{eq:emissions}).

\item[\texttt{max\_queue\_length}] [vehicles] ---
  Worst-case queue across all intersections and steps.  Spikes during
  degradation periods when the equal-phase policy under-serves the dominant
  direction.

\item[\texttt{avg\_queue\_length}] [vehicles] ---
  Network-average queue length over the full simulation window.

\end{description}
